% Relacion de ejercicios de los seminarios 1 y 2: modelo entidad-relacion y
% paso a tablas. Antes esto eran 30 fotografias de la hoja de enunciados con
% los diagramas resueltos a mano; la hoja es del profesorado, asi que el
% enunciado va redactado de nuevo y los diagramas dibujados en tikz.
%
% Notacion del paso a tablas, la misma que se usa en clase:
%   - la clave primaria va subrayada,
%   - CE(n) marca una clave externa y n la relacion a la que apunta,
%   - CC marca una clave candidata.

% Estilo comun de los diagramas E/R. Es el mismo que usa seminario1.tex.
\tikzset{
    entidad/.style={rectangle, draw, minimum width=22mm, minimum height=8mm,
                    align=center, font=\small},
    debil/.style={rectangle, draw, double, double distance=2pt,
                  minimum width=22mm, minimum height=8mm, align=center, font=\small},
    relacion/.style={diamond, draw, aspect=1.9, inner sep=1pt,
                     align=center, font=\footnotesize},
    atributo/.style={ellipse, draw, inner sep=1pt, font=\scriptsize},
    clave/.style={ellipse, draw, inner sep=1pt, font=\scriptsize,
                  label={[font=\tiny]below:\underline{~}}},
    enlace/.style={draw},
    doble/.style={draw, double, double distance=2pt},
    jerarquia/.style={draw, -{Triangle[open, length=3mm, width=4mm]}}
}

\subsection{Tarjetas de crédito, titulares y cuentas}

\textbf{Enunciado.} Se maneja información de tarjetas de crédito, identificadas
por un número y con un tipo; de sus titulares, de los que se conocen el DNI, el
domicilio y el teléfono; y de cuentas corrientes, con código, saldo y fecha de
apertura. Las restricciones son: una persona puede tener más de una tarjeta;
cada tarjeta tiene un único titular; cada tarjeta está asociada a una única
cuenta; se puede cargar más de una tarjeta a una misma cuenta; una cuenta puede
tener asociada varios clientes; y una persona puede tener más de una cuenta.

\begin{figure}[H]
\centering
\begin{tikzpicture}[node distance=14mm and 20mm]
    \node[entidad] (tar) {Tarjeta};
    \node[relacion, right=of tar] (aso) {Asociada};
    \node[entidad, right=of aso] (cc) {Cuenta};
    \node[relacion, below=of aso] (vin) {Vinculada};
    \node[relacion, below left=14mm and 8mm of tar] (tit) {Titular};
    \node[entidad, below=28mm of aso] (per) {Persona};

    \draw (tar) -- (aso);   \draw (aso) -- (cc);
    \draw (cc)  -- (vin);   \draw (vin) -- (per);
    \draw (tar) -- (tit);   \draw (tit) -- (per);

    \node[atributo, above left=6mm and 0mm of tar] (n)  {número};
    \node[atributo, above right=6mm and 0mm of tar] (t) {tipo};
    \draw (tar) -- (n); \draw (tar) -- (t);

    \node[atributo, above=6mm of cc] (cod) {código};
    \node[atributo, right=5mm of cc] (sal) {saldo};
    \node[atributo, below right=4mm and 3mm of cc] (fa) {fecha apertura};
    \draw (cc) -- (cod); \draw (cc) -- (sal); \draw (cc) -- (fa);

    \node[atributo, below left=5mm and 0mm of per] (dni) {DNI};
    \node[atributo, below=5mm of per] (dom) {domicilio};
    \node[atributo, below right=5mm and 0mm of per] (tel) {teléfono};
    \draw (per) -- (dni); \draw (per) -- (dom); \draw (per) -- (tel);
\end{tikzpicture}
\caption{Diagrama E/R del ejercicio 1.}
\end{figure}

\textbf{Paso a tablas.}

\begin{itemize}[leftmargin=*]
    \item \texttt{Persona(\underline{DNI}, domicilio, teléfono)}
    \item \texttt{Cuenta(\underline{código}, saldo, fechaApertura)}
    \item \texttt{Tarjeta(\underline{número}, tipo, DNI\textsuperscript{CE(Persona)},
          código\textsuperscript{CE(Cuenta)})}
    \item \texttt{Vinculada(\underline{código}\textsuperscript{CE(Cuenta)},
          \underline{DNI}\textsuperscript{CE(Persona)})}
\end{itemize}

En \texttt{Tarjeta}, el DNI del titular no admite nulos —toda tarjeta tiene
titular— y el código de cuenta tampoco, porque toda tarjeta está asociada a una.
La relación \emph{vinculada} es de muchos a muchos y por eso necesita tabla
propia.

\subsubsection*{Otra forma}

Si en vez de relacionar la tarjeta con la cuenta se relaciona el titular con la
cuenta y la tarjeta cuelga de esa relación, el resultado es equivalente: la
fusión de las tres relaciones en una deja el mismo esquema relacional. Conviene
comprobarlo, porque es la comprobación que dice si dos diseños distintos
guardan la misma información.

\subsection{Biblioteca: libros, autores, temas y préstamos}

\textbf{Enunciado.} Una biblioteca guarda información de libros, autores, temas,
préstamos y usuarios, con los atributos habituales de cada uno. Restricciones:
un libro puede estar escrito por más de un autor; un autor puede escribir más de
un libro; un libro trata más de un tema; hay muchos libros de cada tema; no
existe más que un ejemplar de cada libro; se registra cuándo un usuario toma
prestado un libro y se elimina el registro cuando lo devuelve, guardando también
la fecha del préstamo; y un usuario no puede tener prestado más de un libro
simultáneamente.

\begin{figure}[H]
\centering
\begin{tikzpicture}[node distance=13mm and 18mm]
    \node[entidad] (lib) {Libro};
    \node[relacion, right=of lib] (esc) {Escribe};
    \node[entidad, right=of esc] (aut) {Autor};
    \node[relacion, below left=13mm and 4mm of lib] (tra) {Trata de};
    \node[entidad, below=26mm of lib] (tem) {Tema};
    \node[relacion, below=of lib.south east, xshift=18mm] (pre) {Préstamo};
    \node[entidad, right=of pre] (usu) {Usuario};

    \draw (lib) -- (esc); \draw (esc) -- (aut);
    \draw (lib) -- (tra); \draw (tra) -- (tem);
    \draw (lib) -- (pre); \draw (pre) -- (usu);

    \node[atributo, above left=6mm and 0mm of lib] (at1) {título}; \draw (lib) -- (at1);
    \node[atributo, above right=6mm and 0mm of lib] (at2) {ISBN}; \draw (lib) -- (at2);
    \node[atributo, above=6mm of aut] (at3) {DNI}; \draw (aut) -- (at3);
    \node[atributo, right=5mm of aut] (at4) {domicilio}; \draw (aut) -- (at4);
    \node[atributo, left=5mm of tem] (at5) {nombre}; \draw (tem) -- (at5);
    \node[atributo, below=5mm of pre] (at6) {fecha}; \draw (pre) -- (at6);
    \node[atributo, below=5mm of usu] (at7) {DNI}; \draw (usu) -- (at7);
\end{tikzpicture}
\caption{Diagrama E/R del ejercicio 2.}
\end{figure}

\textbf{Paso a tablas.}

\begin{itemize}[leftmargin=*]
    \item \texttt{Libro(\underline{ISBN}, título, \dots)}
    \item \texttt{Autor(\underline{DNI}, domicilio, \dots)}
    \item \texttt{Tema(\underline{nombre})}
    \item \texttt{Usuario(\underline{DNI}, nombre)}
    \item \texttt{Escribe(\underline{ISBN}\textsuperscript{CE(Libro)},
          \underline{DNI}\textsuperscript{CE(Autor)})}
    \item \texttt{TrataDe(\underline{ISBN}\textsuperscript{CE(Libro)},
          \underline{nombre}\textsuperscript{CE(Tema)})}
    \item \texttt{Préstamo(\underline{ISBN}\textsuperscript{CE(Libro)}, fecha,
          DNI\textsuperscript{CE(Usuario)})} con \texttt{DNI} como clave
          candidata
\end{itemize}

El DNI del usuario es clave candidata de \texttt{Préstamo} porque un usuario no
puede tener más de un libro a la vez: si pudiera, el par tendría que formar
parte de la clave.

\subsubsection*{Fusión}

Como la relación es de uno a uno por los dos lados y no existe más que un
ejemplar de cada libro, \texttt{Libro} y \texttt{Préstamo} se pueden fusionar:

\begin{center}
\texttt{LibroPréstamo(\underline{ISBN}, título, DNI\textsuperscript{CE(Usuario)}, fecha)}
\end{center}

con \texttt{DNI} como clave candidata. Si en un examen hay que elegir, esta es
la mejor opción, porque es la que menos tablas deja. Deja de valer en cuanto la
clave primaria de las dos entidades no coincide: entonces
\texttt{Préstamo(\underline{fecha, DNI}, \dots)} necesita tabla propia.

\subsection{Biblioteca con ejemplares e histórico de préstamos}

\textbf{Enunciado.} Como el anterior, pero ahora existen varios ejemplares de
cada libro y se guarda el histórico: en el histórico de préstamos se registra una
entrada por cada día que permanece prestada una copia; una misma copia no puede
estar prestada a varios usuarios el mismo día; y un usuario sí puede tener
prestados varios ejemplares al mismo tiempo. Se pide el diagrama y, además,
cómo habría que modificarlo para volver a la restricción de un solo ejemplar por
usuario y día.

\begin{figure}[H]
\centering
\begin{tikzpicture}[node distance=13mm and 18mm]
    \node[entidad] (lib) {Libro};
    \node[relacion, right=of lib] (esc) {Escribe};
    \node[entidad, right=of esc] (aut) {Autor};
    \node[debil, below=of lib] (ej) {Ejemplar};
    \node[relacion, below=of ej] (pre) {Préstamo};
    \node[entidad, right=26mm of pre] (usu) {Usuario};
    \node[relacion, left=22mm of ej] (tra) {Trata de};
    \node[entidad, below=of tra] (tem) {Tema};

    \draw (lib) -- (esc); \draw (esc) -- (aut);
    \draw[doble] (lib) -- (ej);
    \draw[doble] (ej) -- (pre);
    \draw (pre) -- (usu);
    \draw (lib) -- (tra); \draw (tra) -- (tem);

    \node[atributo, above=6mm of lib] (at8) {ISBN}; \draw (lib) -- (at8);
    \node[atributo, right=4mm of ej] (at9) {número}; \draw (ej) -- (at9);
    \node[atributo, left=4mm of pre] (at10) {fecha}; \draw (pre) -- (at10);
    \node[atributo, below=5mm of usu] (at11) {DNI}; \draw (usu) -- (at11);
\end{tikzpicture}
\caption{Diagrama E/R del ejercicio 3. El ejemplar es entidad débil del libro.}
\end{figure}

\textbf{Paso a tablas.}

\begin{itemize}[leftmargin=*]
    \item \texttt{Libro(\underline{ISBN}, título, \dots)}
    \item \texttt{Autor(\underline{DNI}, domicilio, \dots)}
    \item \texttt{Tema(\underline{nombre})}
    \item \texttt{Usuario(\underline{DNI}, nombre)}
    \item \texttt{Escribe(\underline{ISBN}\textsuperscript{CE(Libro)},
          \underline{DNI}\textsuperscript{CE(Autor)})}
    \item \texttt{TrataDe(\underline{ISBN}\textsuperscript{CE(Libro)},
          \underline{nombre}\textsuperscript{CE(Tema)})}
    \item \texttt{Copia(\underline{ISBN}\textsuperscript{CE(Libro)},
          \underline{numCopia})}
    \item \texttt{Préstamo(\underline{DNI}\textsuperscript{CE(Usuario)},
          \underline{fecha}, \underline{ISBN}, \underline{numCopia})}, donde el
          par \texttt{(ISBN, numCopia)} es clave externa a \texttt{Copia}
\end{itemize}

Al ser \texttt{Ejemplar} una entidad débil, su clave la forman la del libro más
un discriminante, y esa clave compuesta viaja entera al préstamo. La clave del
préstamo tiene que incluir la fecha porque se registra una entrada por día.

\textbf{Apartado b.} Para volver a un solo ejemplar prestado por usuario y día,
basta con cambiar la cardinalidad de \emph{Usuario} en la relación de préstamo:
de muchos a uno. En el esquema relacional se traduce en sacar el ejemplar de la
clave y dejar \texttt{Préstamo(\underline{DNI, fecha}, ISBN, numCopia)}.

\subsection{Piezas simples y compuestas}

\textbf{Enunciado.} Una empresa mecánica quiere calcular el precio de las piezas
instaladas en un coche, sabiendo que algunas piezas pueden tener varios
componentes: un motor es batería más ventilador más circuito de arranque. De cada
tipo de pieza se registra un código y una denominación. Hay dos tipos, simple y
compuesta; el precio de una pieza simple consiste en el valor de dicha pieza; el
de una compuesta se corresponde con el precio de montaje sin incluir el de los
tipos de pieza que la componen; para las compuestas se registra el número de
unidades de cada tipo de pieza que la compone; y un tipo de pieza es componente
de un único tipo de pieza compuesta.

\begin{figure}[H]
\centering
\begin{tikzpicture}[node distance=13mm and 16mm]
    \node[entidad] (pz) {Pieza};
    \node[relacion, below=of pz] (esun) {Es un};
    \node[entidad, below left=14mm and 6mm of esun] (sim) {Simple};
    \node[entidad, below right=14mm and 6mm of esun] (com) {Compuesta};
    \node[relacion, right=30mm of pz] (form) {Formada};

    \draw (pz) -- (esun);
    \draw[jerarquia] (sim) -- (esun);
    \draw[jerarquia] (com) -- (esun);
    \draw (pz.east) -- (form.west);
    \draw (form.south) |- (com.east);

    \node[atributo, above left=5mm and 0mm of pz] (at12) {código}; \draw (pz) -- (at12);
    \node[atributo, above right=5mm and 0mm of pz] (at13) {denominación}; \draw (pz) -- (at13);
    \node[atributo, left=4mm of sim] (at14) {precio}; \draw (sim) -- (at14);
    \node[atributo, below=5mm of com] (at15) {precio montaje}; \draw (com) -- (at15);
    \node[atributo, above=5mm of form] (at16) {n.º unidades}; \draw (form) -- (at16);
\end{tikzpicture}
\caption{Diagrama E/R del ejercicio 4. La jerarquía es total y disjunta: una
pieza no puede ser simple y compuesta a la vez.}
\end{figure}

\textbf{Paso a tablas.}

\begin{itemize}[leftmargin=*]
    \item \texttt{TipoPieza(\underline{código}, denominación, stock)}
    \item \texttt{PiezaSimple(\underline{código}\textsuperscript{CE(TipoPieza)},
          precio)}
    \item \texttt{PiezaCompuesta(\underline{código}\textsuperscript{CE(TipoPieza)},
          precioMontaje, precioTotal)}
    \item \texttt{Formada(\underline{codTipo}\textsuperscript{CE(TipoPieza)},
          codCompuesta\textsuperscript{CE(PiezaCompuesta)}, numUnidades)}
\end{itemize}

La clave de \texttt{Formada} es solo el tipo de pieza componente, porque cada
tipo forma parte de una única compuesta. Y una advertencia que conviene tener
presente: \textbf{cuando se usa herencia no tiene sentido fusionar}, porque las
tablas hijas existen precisamente para guardar los atributos que no comparten.

\subsection{Videoclub}

\textbf{Enunciado.} Un videoclub trabaja con películas —título, año de estreno,
actores principales y tema—, cintas —código de cinta y sistema de
reproducción—, préstamos con su fecha y clientes con DNI, nombre, dirección y
teléfono. Un cliente puede alquilar varias cintas el mismo día; puede haber
distintas cintas de la misma película; puede haber películas distintas con el
mismo nombre —versiones—, pero de distinto año; y las películas con el mismo
título son del mismo tema.

\begin{figure}[H]
\centering
\begin{tikzpicture}[node distance=13mm and 18mm]
    \node[entidad] (pel) {Película};
    \node[entidad, right=32mm of pel] (cin) {Cinta};
    \node[relacion, below=of pel] (tra) {Tratan};
    \node[entidad, below=of tra] (tem) {Temas};
    \node[relacion, right=8mm of tra] (act) {Actúan};
    \node[entidad, below=of act] (ac) {Actores};
    \node[relacion, below=of cin] (alq) {Alquile};
    \node[entidad, below=of alq] (cli) {Clientes};

    \draw (pel) -- (cin);
    \draw (pel) -- (tra); \draw (tra) -- (tem);
    \draw (pel) -- (act); \draw (act) -- (ac);
    \draw (cin) -- (alq); \draw[jerarquia] (alq) -- (cli);

    \node[atributo, above left=5mm and 0mm of pel] (at17) {título}; \draw (pel) -- (at17);
    \node[atributo, above right=5mm and 0mm of pel] (at18) {año}; \draw (pel) -- (at18);
    \node[atributo, above=5mm of cin] (at19) {código}; \draw (cin) -- (at19);
    \node[atributo, right=5mm of cin] (at20) {sist. reproducción}; \draw (cin) -- (at20);
    \node[atributo, right=5mm of alq] (at21) {fecha inicio}; \draw (alq) -- (at21);
    \node[atributo, below right=4mm and 2mm of alq] (at22) {fecha fin}; \draw (alq) -- (at22);
    \node[atributo, below=5mm of cli] (at23) {DNI}; \draw (cli) -- (at23);
    \node[atributo, left=4mm of tem] (at24) {código}; \draw (tem) -- (at24);
\end{tikzpicture}
\caption{Diagrama E/R del ejercicio 5.}
\end{figure}

\subsubsection*{Otra versión}

La restricción de que dos películas con el mismo título sean de distinto año se
puede garantizar en el propio diagrama sacando el título a una entidad
\texttt{Título} y el año a otra, unidas por una relación \emph{versiones}: así
la clave de la película es el par y la restricción queda impuesta por el modelo
en vez de por una comprobación externa.

\textbf{Paso a tablas.}

\begin{itemize}[leftmargin=*]
    \item \texttt{Cliente(\underline{DNI}, nombre, dirección, teléfono)}
    \item \texttt{Tema(\underline{numTema})}
    \item \texttt{Título(\underline{numTítulo})}
    \item \texttt{Año(\underline{numAño})}
    \item \texttt{Cinta(\underline{numTítulo}, numAño, numCinta, tipo)}
    \item \texttt{Película(\underline{numTítulo}\textsuperscript{CE(Título)},
          \underline{año}\textsuperscript{CE(Año)}, reparto)}
    \item \texttt{Tratan(\underline{numTítulo}\textsuperscript{CE(Título)},
          numTema\textsuperscript{CE(Tema)})}
    \item \texttt{Préstamo(\underline{DNI}\textsuperscript{CE(Cliente)},
          \underline{fecha}, \underline{numTítulo}, \underline{año},
          \underline{numCinta})}
\end{itemize}

\texttt{Título} y \texttt{Tratan} se pueden fusionar, porque la relación es de
uno a uno y obligatoria por el lado del título.

\subsection{Modelo de factura}

\textbf{Enunciado.} Realizar el diagrama E/R que permita generar la información
de una factura: emisor y destinatario con nombre, dirección y CIF; número y
fecha de factura; y una línea por artículo con cantidad, descripción, precio
unitario, IVA y el importe resultante.

\begin{figure}[H]
\centering
\begin{tikzpicture}[node distance=13mm and 18mm]
    \node[entidad] (emi) {Emisor};
    \node[relacion, right=of emi] (fac) {Factura};
    \node[entidad, right=of fac] (des) {Destinatario};
    \node[relacion, below=20mm of fac] (tie) {Tiene};
    \node[entidad, below=of tie] (art) {Artículos};

    \draw (emi) -- (fac); \draw (fac) -- (des);
    \draw[jerarquia] (tie) -- (fac);
    \draw (tie) -- (art);

    \node[atributo, above=5mm of fac] (at25) {fecha}; \draw (fac) -- (at25);
    \node[atributo, below left=6mm and 0mm of emi] (at26) {CIF}; \draw (emi) -- (at26);
    \node[atributo, above=6mm of des] (at27) {CIF}; \draw (des) -- (at27);
    \node[atributo, right=5mm of tie] (at28) {cantidad}; \draw (tie) -- (at28);
    \node[atributo, below right=4mm and 2mm of tie] (at29) {importe}; \draw (tie) -- (at29);
    \node[atributo, below=5mm of art] (at30) {precio}; \draw (art) -- (at30);
    \node[atributo, right=5mm of art] (at31) {IVA}; \draw (art) -- (at31);
\end{tikzpicture}
\caption{Diagrama E/R del ejercicio 6, primera solución.}
\end{figure}

Esta solución solo distingue emisores y receptores. Cuando las mismas empresas
son unas veces emisor y otras cliente, duplica información, y esa es la razón de
la versión optimizada: una única entidad \texttt{Empresa} y dos relaciones,
\emph{emite} y \emph{recibe}, hacia ella.

\textbf{Paso a tablas} de la versión optimizada:

\begin{itemize}[leftmargin=*]
    \item \texttt{Empresa(\underline{NIF}, nombre, dirección)}
    \item \texttt{Artículo(\underline{numArt}, descripción, precioVenta, IVA)}
    \item \texttt{Factura(\underline{numFactura}, fecha, baseImponible, IVA,
          total)}
    \item \texttt{Emite(\underline{numFactura}\textsuperscript{CE(Factura)},
          NIF\textsuperscript{CE(Empresa)})}
    \item \texttt{Recibe(\underline{numFactura}\textsuperscript{CE(Factura)},
          NIF\textsuperscript{CE(Empresa)})}
    \item \texttt{Línea(\underline{numLínea},
          \underline{numFactura}\textsuperscript{CE(Factura)})}
    \item \texttt{Incluye(\underline{numLínea}, \underline{numFactura},
          numArt\textsuperscript{CE(Artículo)}, cantidad)}
\end{itemize}

Los importes de la factura —base imponible, IVA y total— son atributos
calculados: dependen del precio y del IVA de cada artículo y de la cantidad.
Guardarlos es una decisión de diseño, no una necesidad del modelo.

\subsubsection*{Fusiones}

\texttt{Factura}, \texttt{Emite} y \texttt{Recibe} se pueden fusionar en una
sola tabla con los dos NIF, y \texttt{Línea} con \texttt{Incluye}. Esta segunda
fusión es obligatoria, no opcional, porque la relación es obligatoria por el
lado de la línea: una línea de factura sin artículo no existe.

\subsection{Tienda de productos informáticos}

\textbf{Enunciado.} En la base de datos de una tienda de productos informáticos,
los tipos de producto se registran con un número de referencia, un fabricante y
un precio de venta al público. De los artículos estrella —impresoras,
ordenadores personales y portátiles— se registran además sus características
específicas: color, resolución y tipo en las impresoras; procesador, velocidad,
RAM y capacidad de disco en los PC; y lo mismo más el tamaño de pantalla en los
portátiles.

\begin{figure}[H]
\centering
\begin{tikzpicture}[node distance=14mm and 10mm]
    \node[entidad] (pro) {Producto};
    \node[relacion, below=of pro] (esun) {Es un};
    \node[entidad, below left=16mm and 14mm of esun] (imp) {Impresora};
    \node[entidad, below=16mm of esun] (pc) {PC};
    \node[entidad, below right=16mm and 14mm of esun] (por) {Portátil};

    \draw (pro) -- (esun);
    \draw[jerarquia] (imp) -- (esun);
    \draw[jerarquia] (pc)  -- (esun);
    \draw[jerarquia] (por) -- (esun);

    \node[atributo, above=5mm of pro] (at32) {n.º referencia}; \draw (pro) -- (at32);
    \node[atributo, left=5mm of pro] (at33) {fabricante}; \draw (pro) -- (at33);
    \node[atributo, right=5mm of pro] (at34) {PVP}; \draw (pro) -- (at34);
    \node[atributo, below=5mm of imp] (at35) {color, resolución, tipo}; \draw (imp) -- (at35);
    \node[atributo, below=5mm of pc] (at36) {procesador, RAM\dots}; \draw (pc) -- (at36);
    \node[atributo, below=5mm of por] (at37) {\dots\ tamaño pantalla}; \draw (por) -- (at37);
\end{tikzpicture}
\caption{Diagrama E/R del ejercicio 7. Jerarquía parcial: hay productos que no
son de ninguno de los tres tipos.}
\end{figure}

La jerarquía es \emph{parcial} porque la tienda vende más cosas que impresoras,
PC y portátiles, y \emph{disjunta} porque un producto no puede ser de dos tipos
a la vez. En el paso a tablas, cada hija lleva la clave del padre como clave
primaria y clave externa a la vez.

\subsection{Cine, salas y proyecciones}

\textbf{Enunciado.} Un cine tiene varias salas donde se proyectan películas en
una hora y fecha determinadas. Las salas se componen de asientos identificados
por fila y número, y hay que saber si en cada proyección un asiento está libre u
ocupado. El cine compra una o varias copias de una película para poder
proyectarla en distintas salas. Restricciones: cada sala tiene un número
determinado de asientos; una película se identifica por un código y se describe
por título y duración; una película puede tener varias copias; una proyección se
identifica por la sala, la copia, la fecha y la hora; en una sala no pueden
proyectarse el mismo día y a la misma hora dos cintas distintas; cada entrada
identifica la sala, la película, la fecha, la hora de comienzo de la proyección,
la fila y el asiento; y no pueden venderse entradas que supongan la ocupación
del mismo asiento a la vez.

\begin{figure}[H]
\centering
\begin{tikzpicture}[node distance=13mm and 18mm]
    \node[entidad] (sala) {Sala};
    \node[relacion, right=26mm of sala] (proy) {Proyecta};
    \node[entidad, right=of proy] (cop) {Copias};
    \node[entidad, below=of cop] (pel) {Película};
    \node[debil, below=of sala] (asi) {Asiento};
    \node[relacion, below right=16mm and 4mm of asi] (vta) {Venta de entradas};

    \draw (sala) -- (proy); \draw (proy) -- (cop);
    \draw[doble] (cop) -- (pel);
    \draw[doble] (sala) -- (asi);
    \draw[doble] (asi) -- (vta);
    \draw (vta) -- (proy);

    \node[atributo, above left=5mm and 0mm of sala] (at38) {núm. sala}; \draw (sala) -- (at38);
    \node[atributo, above right=5mm and 0mm of sala] (at39) {n.º asientos}; \draw (sala) -- (at39);
    \node[atributo, below left=5mm and 0mm of proy] (at40) {fecha}; \draw (proy) -- (at40);
    \node[atributo, below right=5mm and 0mm of proy] (at41) {hora}; \draw (proy) -- (at41);
    \node[atributo, above=5mm of cop] (at42) {n.º copia}; \draw (cop) -- (at42);
    \node[atributo, right=5mm of pel] (at43) {título, duración, código}; \draw (pel) -- (at43);
    \node[atributo, below left=5mm and 0mm of asi] (at44) {fila}; \draw (asi) -- (at44);
    \node[atributo, below right=9mm and 4mm of asi] (at45) {columna}; \draw (asi) -- (at45);
\end{tikzpicture}
\caption{Diagrama E/R del ejercicio 8.}
\end{figure}

\textbf{Apartado b.} El diseño no garantiza por sí solo que no se solapen en el
tiempo dos proyecciones distintas en una misma sala: la clave de la proyección
incluye la copia, así que dos copias distintas pueden compartir sala, fecha y
hora. Para impedirlo hay que sacar la copia de la clave y dejarla como atributo,
de forma que la clave sea \texttt{(sala, fecha, hora)}. Con eso el propio modelo
impone la restricción, en vez de dejarla a un disparador.

\subsection{Registro de llamadas}

\textbf{Enunciado.} Expresar mediante un diagrama E/R el registro de llamadas
entre dos teléfonos, con fecha y hora de inicio y de finalización. Un teléfono
se identifica por un número y hay dos tipos: fijo y móvil. No se contemplan
llamadas en las que participen más de dos teléfonos.

\begin{figure}[H]
\centering
\begin{tikzpicture}[node distance=14mm and 26mm]
    \node[entidad] (tel) {Teléfono};
    \node[relacion, right=of tel] (lla) {Llamada};

    \draw (tel.north) to[out=60, in=150] node[above, font=\scriptsize] {llama} (lla.north west);
    \draw (tel.south) to[out=-60, in=-150] node[below, font=\scriptsize] {recibe} (lla.south west);

    \node[atributo, above=6mm of lla] (at46) {fecha}; \draw (lla) -- (at46);
    \node[atributo, right=5mm of lla] (at47) {hora inicio}; \draw (lla) -- (at47);
    \node[atributo, below=6mm of lla] (at48) {hora fin}; \draw (lla) -- (at48);
    \node[atributo, left=5mm of tel] (at49) {número}; \draw (tel) -- (at49);
\end{tikzpicture}
\caption{Diagrama E/R del ejercicio 9. La relación es reflexiva: el mismo
conjunto de entidades participa con dos papeles distintos.}
\end{figure}

Lo que hace especial a este ejercicio es que la relación es \emph{reflexiva} y
hay que etiquetar los dos papeles, porque «llama» y «recibe» no son
intercambiables. En el paso a tablas la llamada necesita los dos números como
claves externas al mismo teléfono, con nombres distintos.

\subsection{Recetas de cocina}

\textbf{Enunciado.} Una receta se describe mediante una serie de ingredientes y
de pasos de ejecución. Las recetas tienen código, nombre, tipo —primero,
segundo o postre— y dificultad —alto, medio o bajo—. Los ingredientes tienen
código, nombre, tipo —grano, polvo, troceado u otro— y precio. Realizar el
diagrama para almacenar las recetas completas.

\begin{figure}[H]
\centering
\begin{tikzpicture}[node distance=14mm and 20mm]
    \node[entidad] (rec) {Receta};
    \node[entidad, right=36mm of rec] (ing) {Ingrediente};
    \node[relacion, below=of rec] (cre) {Crean};
    \node[relacion, below=of ing] (usa) {Usan};
    \node[debil, below right=14mm and 2mm of cre] (pas) {Pasos};

    \draw (rec) -- (cre); \draw[doble] (cre) -- (pas);
    \draw (ing) -- (usa); \draw (usa) -- (pas);

    \node[atributo, above left=5mm and 0mm of rec] (at50) {CODR}; \draw (rec) -- (at50);
    \node[atributo, above right=5mm and 0mm of rec] (at51) {tipo, dificultad}; \draw (rec) -- (at51);
    \node[atributo, above=5mm of ing] (at52) {CODI}; \draw (ing) -- (at52);
    \node[atributo, right=5mm of ing] (at53) {tipo, precio}; \draw (ing) -- (at53);
    \node[atributo, left=5mm of cre] (at54) {n.º pasos}; \draw (cre) -- (at54);
    \node[atributo, right=5mm of usa] (at55) {cantidad}; \draw (usa) -- (at55);
    \node[atributo, below=5mm of pas] (at56) {descripción}; \draw (pas) -- (at56);
\end{tikzpicture}
\caption{Diagrama E/R del ejercicio 10.}
\end{figure}

El paso es entidad débil de la receta: fuera de ella no tiene sentido, y su
identificador es el número de orden dentro de la receta. La cantidad va en la
relación \emph{usan} y no en el ingrediente, porque depende del par.

\subsection{Artículos científicos en revistas}

\textbf{Enunciado.} Se quiere gestionar la publicación de artículos científicos.
Una revista se identifica por su ISSN y tiene nombre y editorial. Durante un año
publica uno o varios números, que recogen los artículos aceptados; el número
identifica a cada publicación de la revista y se recoge también el año. Cada
número contiene varios artículos. Un artículo se identifica por el título y por
el número de la revista en que se ha publicado, y se almacena la página de
inicio y de fin en ese número. Cada autor se identifica por un código y se
caracteriza por su nombre y nacionalidad; un artículo puede estar escrito por
varios autores y un autor puede escribir varios artículos; y un artículo puede
hacer referencia a otros artículos y puede ser citado en otros.

\begin{figure}[H]
\centering
\begin{tikzpicture}[node distance=14mm and 20mm]
    \node[entidad] (rev) {Revista};
    \node[debil, below=of rev] (num) {Números};
    \node[debil, below=of num] (art) {Artículo};
    \node[relacion, right=of art] (esc) {Escribe};
    \node[entidad, right=of esc] (aut) {Autor};
    \node[relacion, left=18mm of art] (ref) {Referencia};

    \draw[doble] (rev) -- (num);
    \draw[doble] (num) -- (art);
    \draw (art) -- (esc); \draw (esc) -- (aut);
    \draw (art.west) to[out=170, in=30] node[above, font=\scriptsize, pos=0.3] {cita} (ref.north east);
    \draw (art.south west) to[out=200, in=-30] node[below, font=\scriptsize, pos=0.3] {citado} (ref.south east);

    \node[atributo, above=5mm of rev] (at57) {ISSN}; \draw (rev) -- (at57);
    \node[atributo, right=5mm of rev] (at58) {nombre, editorial}; \draw (rev) -- (at58);
    \node[atributo, right=5mm of num] (at59) {n.º, año}; \draw (num) -- (at59);
    \node[atributo, below=5mm of art] (at60) {título, pág. ini, pág. fin}; \draw (art) -- (at60);
    \node[atributo, below=5mm of aut] (at61) {cód., nombre, nacionalidad}; \draw (aut) -- (at61);
\end{tikzpicture}
\caption{Diagrama E/R del ejercicio 11.}
\end{figure}

\textbf{Paso a tablas.}

\begin{itemize}[leftmargin=*]
    \item \texttt{Revista(\underline{ISSN}, nombre, editorial)}
    \item \texttt{Autor(\underline{numAutor}, nombre, nacionalidad)}
    \item \texttt{Número(\underline{numRevista}, \underline{año},
          ISSN\textsuperscript{CE(Revista)})}
    \item \texttt{Artículo(\underline{ISSN}, \underline{num}, \underline{año},
          \underline{título}, págIni, págFin)}
    \item \texttt{Escribe(\underline{numAutor}\textsuperscript{CE(Autor)},
          \underline{ISSN}, \underline{num}, \underline{año},
          \underline{título})}
    \item \texttt{Referencia(\underline{ISSN\_c}, \underline{num\_c},
          \underline{año\_c}, \underline{título\_c}, \dots\ el citado)}
\end{itemize}

En \texttt{Referencia}, el sufijo distingue el artículo que cita del citado: la
tabla lleva dos veces la misma clave compuesta con nombres distintos, igual que
en el ejercicio de las llamadas. Y no se puede fusionar nada aquí, porque todas
las relaciones son de muchos a muchos.

\subsection{Reservas de vuelos}

\textbf{Enunciado.} Se organiza la gestión de reservas de vuelos. En la tarjeta
de embarque figuran el pasajero, la fecha y hora de emisión, el asiento que le
corresponde, el avión, la fecha y hora de salida y el destino —ciudad de partida
y de llegada—. Restricciones: hay aviones diferentes cuyos números de asiento
pueden coincidir; cada avión tiene una capacidad máxima; un destino aéreo viene
identificado por un número y puede incluir varios trayectos especificados por la
ciudad de salida y la de llegada; un vuelo es la realización de un trayecto
mediante un avión, con una fecha y hora de partida determinadas; un avión puede
participar en diferentes vuelos; una tarjeta de embarque corresponde a un
asiento concreto, de un avión concreto, en un vuelo concreto y para un pasajero
concreto; se emiten varias tarjetas de embarque para cada vuelo; y en un vuelo,
un asiento no puede estar ocupado por más de un pasajero.

\begin{figure}[H]
\centering
\begin{tikzpicture}[node distance=13mm and 16mm]
    \node[entidad] (tar) {Tarjeta de\\embarque};
    \node[relacion, right=of tar] (cor) {Corresponde};
    \node[debil, right=of cor] (asi) {Asiento};
    \node[relacion, below=of asi] (tie) {Tiene};
    \node[entidad, below=of tie] (avi) {Avión};
    \node[relacion, left=of avi] (rea) {Realiza};
    \node[entidad, left=of rea] (vue) {Vuelo};
    \node[relacion, below=of vue] (inc) {Incluye};
    \node[entidad, right=of inc] (des) {Destino aéreo};
    \node[entidad, below=of tar] (pas) {Pasajero};

    \draw (tar) -- (cor); \draw (cor) -- (asi);
    \draw[doble] (asi) -- (tie); \draw[doble] (tie) -- (avi);
    \draw (avi) -- (rea); \draw (rea) -- (vue);
    \draw (vue) -- (inc); \draw (inc) -- (des);
    \draw (tar) -- (pas);
    \draw (vue.north) |- (tar.west);

    \node[atributo, above=5mm of tar] (at62) {cód., fecha y hora emisión}; \draw (tar) -- (at62);
    \node[atributo, below=5mm of pas] (at63) {DNI}; \draw (pas) -- (at63);
    \node[atributo, right=5mm of avi] (at64) {cód., modelo, capacidad}; \draw (avi) -- (at64);
    \node[atributo, above=5mm of vue] (at65) {fecha, hora salida}; \draw (vue) -- (at65);
    \node[atributo, below=5mm of des] (at66) {ciudad salida, ciudad llegada}; \draw (des) -- (at66);
\end{tikzpicture}
\caption{Diagrama E/R del ejercicio 12.}
\end{figure}

\textbf{Apartado b.} Un vuelo se identifica por el avión, el destino, la ciudad
de llegada y la ciudad de salida. Esa información se recoge en las relaciones
\emph{realiza} e \emph{incluye}, y por tanto la clave del vuelo la forman las
claves externas a \texttt{Avión} y a \texttt{Destino} más la fecha y la hora.

La capacidad máxima del avión está como atributo, así que \textbf{sí} se pueden
emitir más tarjetas de embarque que asientos disponibles: nada en el modelo lo
impide. Para impedirlo habría que contar las tarjetas emitidas del vuelo y
compararlas con la capacidad, y eso es una restricción que el diagrama E/R no
expresa.

\subsection{Programación de canales de televisión}

\textbf{Enunciado.} Se modela la programación que ofrecen los canales de TV.
Presentadores con DNI, nombre y especialidad; canales con nombre, empresa y
tipo; y programas con código, nombre, duración y tipo. Existen tres tipos de
programa: películas, de las que hay que conocer título, tema y año; concursos; e
informativos. Un programa solo puede emitirse por un canal un determinado día a
una determinada hora. Los concursos solo pueden ser presentados por un
presentador, pero un presentador puede serlo de varios concursos. Los
informativos pueden ser presentados por varios presentadores y un presentador
puede serlo de varios también.

\begin{figure}[H]
\centering
\begin{tikzpicture}[node distance=13mm and 16mm]
    \node[entidad] (pre) {Presentadores};
    \node[relacion, right=26mm of pre] (emi) {Emite};
    \node[entidad, right=of emi] (pro) {Programas};
    \node[relacion, below=of pro] (esun) {Es un};
    \node[entidad, below left=15mm and 10mm of esun] (pel) {Película};
    \node[entidad, below=15mm of esun] (con) {Concursos};
    \node[entidad, below right=15mm and 10mm of esun] (inf) {Informativos};
    \node[entidad, below=30mm of pre] (can) {Canales};
    \node[relacion, below=18mm of can] (pst) {Presentan};

    \draw (pro) -- (emi); \draw (emi) -- (can);
    \draw (pro) -- (esun);
    \draw[jerarquia] (pel) -- (esun);
    \draw[jerarquia] (con) -- (esun);
    \draw[jerarquia] (inf) -- (esun);
    \draw (pre) -- (pst);
    \draw (pst) -- (con); \draw (pst) -- (inf);

    \node[atributo, above=5mm of pre] (at67) {DNI, nombre, especialidad}; \draw (pre) -- (at67);
    \node[atributo, above=5mm of pro] (at68) {cód., nombre, duración, tipo}; \draw (pro) -- (at68);
    \node[atributo, below left=5mm and 0mm of emi] (at69) {día}; \draw (emi) -- (at69);
    \node[atributo, below right=5mm and 0mm of emi] (at70) {hora}; \draw (emi) -- (at70);
    \node[atributo, left=5mm of can] (at71) {nombre, empresa, tipo}; \draw (can) -- (at71);
    \node[atributo, below left=7mm and 0mm of pel] (at72) {título, tema, año}; \draw (pel) -- (at72);
\end{tikzpicture}
\caption{Diagrama E/R del ejercicio 13.}
\end{figure}

La jerarquía vuelve a ser total y disjunta. La diferencia entre concursos e
informativos está en la cardinalidad de \emph{presentan}: uno a muchos para los
concursos, muchos a muchos para los informativos, y por eso no se puede
dibujar una sola relación desde el programa genérico.

\subsection{Comisarías de policía}

\textbf{Enunciado.} Se diseña la base de datos de las comisarías de una ciudad,
que registra delincuentes —identificados por su DNI, con nombre y apellidos,
fecha de nacimiento, domicilio y nacionalidad—, delitos —con un código, y de los
que hay que registrar nombre, tipo y penas máxima y mínima—, agentes
—identificados por su número de placa, con DNI, nombre y apellidos, domicilio y
fecha de ingreso en el cuerpo— y comisarías —con nombre y localización—.
Restricciones: existen dos tipos de agente, oficiales y policías, y de los
oficiales se conoce la especialización y los años de experiencia; un oficial es
responsable de un grupo de policías, pero cada policía solo depende de un
oficial; cada comisaría la dirige un oficial; cuando se detiene a un delincuente
hay que registrar el lugar, la fecha y la hora de la detención y qué agentes han
participado; y de los delitos cometidos hay que conocer el lugar, la fecha y la
hora, así como los delincuentes que han participado.

\begin{figure}[H]
\centering
\begin{tikzpicture}[node distance=13mm and 15mm]
    \node[entidad] (del) {Delincuente};
    \node[relacion, right=16mm of del] (com) {Comete};
    \node[entidad, right=16mm of com] (dl) {Delitos};

    \node[relacion, below=16mm of del] (det) {Detiene};
    \node[entidad, below=16mm of det] (ag) {Agente};
    \node[relacion, right=16mm of ag] (par) {Pertenece};
    \node[entidad, right=16mm of par] (cm) {Comisaría};

    \node[relacion, below=16mm of ag] (esun) {Es un};
    \node[entidad, below left=15mm and 2mm of esun] (ofi) {Oficial};
    \node[entidad, below right=15mm and 2mm of esun] (pol) {Policía};
    \node[relacion, below=15mm of ofi] (resp) {Responsable};
    \node[relacion, below=15mm of cm] (dir) {Dirige};

    \draw (del) -- (com); \draw (com) -- (dl);
    \draw (del) -- (det); \draw (det) -- (ag);
    \draw (dl.south) to[out=-90, in=0] (det.east);
    \draw (ag) -- (par); \draw (par) -- (cm);
    \draw (cm) -- (dir); \draw (dir.south) -- ++(0,-26mm) -| (ofi.south);
    \draw (ag) -- (esun);
    \draw[jerarquia] (ofi) -- (esun);
    \draw[jerarquia] (pol) -- (esun);
    \draw (ofi) -- (resp); \draw (resp.east) -| (pol.south);

    \node[atributo, above=5mm of del] (ap1) {DNI, nombre, nacionalidad}; \draw (del) -- (ap1);
    \node[atributo, above=5mm of dl] (ap2) {código, tipo, penas}; \draw (dl) -- (ap2);
    \node[atributo, below=5mm of com] (ap3) {lugar, fecha, hora}; \draw (com) -- (ap3);
    \node[atributo, left=5mm of det] (ap4) {lugar, fecha, hora}; \draw (det) -- (ap4);
    \node[atributo, left=5mm of ag] (ap5) {n.º placa, DNI}; \draw (ag) -- (ap5);
    \node[atributo, above=5mm of cm] (ap6) {nombre, localización}; \draw (cm) -- (ap6);
    \node[atributo, left=5mm of ofi] (ap7) {especialización}; \draw (ofi) -- (ap7);
    \end{tikzpicture}
\caption{Diagrama E/R del ejercicio 14, que es el típico de examen.}
\end{figure}

\textbf{Paso a tablas.}

\begin{itemize}[leftmargin=*]
    \item \texttt{Delincuente(\underline{DNI\_D}, \dots)}
    \item \texttt{Delito(\underline{codDel}, \dots)}
    \item \texttt{Agente(\underline{DNI\_A}, \dots)}
    \item \texttt{Oficial(\underline{DNI\_O}\textsuperscript{CE(Agente)},
          especialización, años)}
    \item \texttt{Policía(\underline{DNI\_P}\textsuperscript{CE(Agente)})}
    \item \texttt{Comisaría(\underline{codCom}, \dots)}
    \item \texttt{Comete(\underline{DNI\_D}\textsuperscript{CE(Delincuente)},
          \underline{codDel}\textsuperscript{CE(Delito)}, \underline{fecha},
          \underline{hora})}
    \item \texttt{Detiene(\underline{DNI\_A}\textsuperscript{CE(Agente)},
          \underline{DNI\_D}, \underline{codDel}, \underline{fecha},
          \underline{hora}, fechaDet, horaDet)}
    \item \texttt{Responsable(\underline{DNI\_P}\textsuperscript{CE(Policía)},
          DNI\_O\textsuperscript{CE(Oficial)})}
    \item \texttt{Dirige(\underline{codCom}\textsuperscript{CE(Comisaría)},
          DNI\_O\textsuperscript{CE(Oficial)})} con \texttt{DNI\_O} como clave
          candidata
\end{itemize}

\subsubsection*{Fusiones}

\texttt{Policía} y \texttt{Responsable} se pueden fusionar, porque cada policía
depende de un solo oficial; y \texttt{Comisaría} con \texttt{Dirige}, por lo
mismo. En el segundo caso el DNI del oficial se queda además como clave
candidata, porque cada oficial dirige como mucho una comisaría.

Dos cosas que el enunciado no cierra y hay que decidir. La primera es si el lugar
del delito se pone como atributo de la relación o como entidad propia: siendo un
dato que se repite en muchos delitos, una entidad ahorra espacio, y complica el
esquema sin necesidad si nadie consulta por él. La segunda es cuándo un
delincuente pasa a serlo; si se registra a alguien antes de imputarle un delito,
su participación en \emph{comete} no puede ser obligatoria.

\subsection{Empresa de alquiler de inmuebles}

\textbf{Enunciado.} De cada inmueble se desea conocer su número de registro
—único—, el tipo, los metros cuadrados y la dirección. Hay dos tipos: pisos y
locales comerciales. De los pisos interesan el número de habitaciones, el de
baños, la planta y, para cada habitación, su tipo y tamaño. De los locales, si
está a estrenar o no y qué tipo de negocios o instalaciones ha tenido
previamente, con tipo de instalación, negocio y costos. De los clientes, el DNI,
el nombre, el teléfono y el número de cuenta. Restricciones: los clientes
alquilan los inmuebles por periodos de tiempo, y durante ese periodo el inmueble
solo puede estar alquilado a un inquilino; el precio del alquiler puede variar
por periodos o inquilinos; y los locales pueden haber tenido instalados
diferentes negocios en diferentes fechas.

\begin{figure}[H]
\centering
\begin{tikzpicture}[node distance=13mm and 16mm]
    \node[entidad] (inm) {Inmueble};
    \node[relacion, right=of inm] (alq) {Alquilan};
    \node[entidad, right=of alq] (cli) {Clientes};
    \node[relacion, below=of inm] (esun) {Es un};
    \node[entidad, below left=15mm and 2mm of esun] (pis) {Pisos};
    \node[entidad, below right=15mm and 2mm of esun] (loc) {Locales};
    \node[debil, below=of pis] (hab) {Habitación};
    \node[relacion, right=of loc] (ins) {Instalan};
    \node[entidad, right=of ins] (neg) {Negocio};

    \draw (inm) -- (alq); \draw (alq) -- (cli);
    \draw (inm) -- (esun);
    \draw[jerarquia] (pis) -- (esun);
    \draw[jerarquia] (loc) -- (esun);
    \draw[doble] (pis) -- (hab);
    \draw (loc) -- (ins); \draw (ins) -- (neg);

    \node[atributo, above left=5mm and 0mm of inm] (at80) {n.º registro}; \draw (inm) -- (at80);
    \node[atributo, above=5mm of alq] (at81) {periodo, precio}; \draw (alq) -- (at81);
    \node[atributo, below=5mm of cli] (at82) {DNI, teléfono, cuenta}; \draw (cli) -- (at82);
    \node[atributo, below=5mm of hab] (at83) {tipo, tamaño}; \draw (hab) -- (at83);
    \node[atributo, left=5mm of loc] (at84) {a estrenar}; \draw (loc) -- (at84);
    \node[atributo, above=5mm of ins] (at85) {fecha}; \draw (ins) -- (at85);
\end{tikzpicture}
\caption{Diagrama E/R del ejercicio 15.}
\end{figure}

La jerarquía aquí es \emph{parcial y disjunta}: el enunciado dice que se
consideran inmuebles de dos tipos, así que un inmueble no puede ser piso y local
a la vez, pero nada obliga a que todo inmueble sea uno de los dos. El periodo va
como atributo de la relación de alquiler y forma parte de su clave: sin él, un
mismo inmueble no podría alquilarse dos veces al mismo cliente.

\subsection{Objetos astronómicos}

\textbf{Enunciado.} Cada objeto astronómico se identifica por un código y
registra la fecha y el observatorio donde se hizo el descubrimiento. Se
clasifican en dos: de baja emisión de luz —planetas y satélites— y de alta
emisión —estrellas—. De los planetas se almacena el tipo, de los satélites el
tipo de satélite, y de las estrellas el tipo y el subtipo. Además: un satélite
gira alrededor de un planeta y se almacena a cuántos años luz se encuentran entre
sí; alrededor de un planeta pueden girar diferentes satélites; un planeta, junto
con sus satélites, gira alrededor de una estrella, y se almacena a cuántos años
luz distan entre sí; alrededor de una estrella pueden girar diferentes planetas;
un grupo de estrellas forma una constelación, cada estrella puede estar en varias
constelaciones, y de la constelación interesan el código y el nombre.

\begin{figure}[H]
\centering
\begin{tikzpicture}[node distance=12mm and 14mm]
    \node[entidad] (obj) {Objeto\\astronómico};
    \node[relacion, below=of obj] (e1) {Es un};
    \node[entidad, below left=14mm and 4mm of e1] (baja) {Baja emisión};
    \node[entidad, below right=14mm and 4mm of e1] (alta) {Alta emisión};
    \node[relacion, below=16mm of baja] (e2) {Es un};
    \node[entidad, below left=13mm and 0mm of e2] (pla) {Planeta};
    \node[entidad, below right=13mm and 0mm of e2] (sat) {Satélite};
    \node[entidad, right=18mm of alta] (est) {Estrella};
    \node[relacion, below=14mm of est] (for) {Forman};
    \node[entidad, below=14mm of for] (con) {Constelación};
    \node[relacion, below=14mm of alta] (gir) {Giran};
    \node[relacion, below=16mm of sat] (alr) {Gira alrededor};

    \draw (obj) -- (e1);
    \draw[jerarquia] (baja) -- (e1);
    \draw[jerarquia] (alta) -- (e1);
    \draw (baja) -- (e2);
    \draw[jerarquia] (pla) -- (e2);
    \draw[jerarquia] (sat) -- (e2);
    \draw (alta) -- (est);
    \draw (est) -- (for); \draw (for) -- (con);
    \draw (est.south) to[out=-90, in=0] (gir.east);
    \draw (gir.west) -| (pla.north);
    \draw (sat) -- (alr); \draw (alr.west) -| (pla.south);

    \node[atributo, above=5mm of obj] (as1) {código, fecha, observatorio}; \draw (obj) -- (as1);
    \node[atributo, above=5mm of est] (as2) {tipo, subtipo}; \draw (est) -- (as2);
    \node[atributo, below=5mm of gir] (as3) {años luz}; \draw (gir) -- (as3);
    \node[atributo, right=5mm of alr] (as4) {años luz}; \draw (alr) -- (as4);
    \node[atributo, below=5mm of con] (as5) {código, nombre}; \draw (con) -- (as5);
    \end{tikzpicture}
\caption{Diagrama E/R del ejercicio 16, con dos niveles de jerarquía.}
\end{figure}

Es el ejercicio que más jerarquías encadena: la primera separa por emisión de
luz y la segunda, dentro de la baja, separa planetas de satélites. Las dos son
totales y disjuntas. Las dos relaciones de giro llevan el mismo atributo con el
mismo significado y aun así no se pueden unir, porque relacionan pares de
entidades distintos.
